\section{Generalization Beyond Euclidean Cost}\label{sec:exp3-cost}

The Euclidean metric used throughout the main experiments is a
modelling convenience: real deployments rarely offer unobstructed
straight-line travel, and real fleets are rarely homogeneous.
Experiment~3 asks whether the $p/d^2$ recommendation survives when
the cost function is perturbed in three realistic ways. For each
variant we run $500$ trials at the default configuration
($n = 30$, $R = 3$, $E = 14$, $L = 10$). The three variants are: a
\emph{Manhattan} metric
$d(p, q) = |p_x - q_x| + |p_y - q_y|$, which models a rectilinear
warehouse or an urban grid in which axis-aligned motion is the only
option; a \emph{heterogeneous-speed} fleet in which each robot $r$
draws a speed multiplier $v_r \sim \mathrm{Uniform}(0.7, 1.3)$ and
its effective travel cost is $d_{\mathrm{Eucl}}(p, q) / v_r$; and an
\emph{obstacle} variant in which three fixed circular obstacles of
radius $0.5$ sit in the interior of the square and each segment
that crosses an obstacle incurs a fixed detour penalty of
$2\,r_{\mathrm{obs}}$. The obstacle radius is calibrated so that
roughly $27\%$ of site-to-site segments cross at least one obstacle.
Figure~\ref{fig:cost-sensitivity} reports the mean Finder
Competitive Ratio for M4 ($p/d$) and M4* ($p/d^2$) under each cost
model.

\begin{figure}[ht]
\centering
\includegraphics[width=0.88\textwidth]{fig_cost_sensitivity.png}
\caption{Mean FCR under four cost models. M4* ($p/d^2$) wins in all
four settings. The margin is largest under the obstacle penalty
($+9.5\%$) and smallest under the Manhattan metric ($+1.5\%$).}
\label{fig:cost-sensitivity}
\end{figure}

The headline result is that \textbf{$p/d^2$ wins under all four cost
models}: mean FCR values are $2.84$ versus $3.10$ under Euclidean
($+8.6\%$ margin in favour of M4*), $3.09$ versus $3.14$ under
Manhattan ($+1.5\%$), $2.89$ versus $3.10$ under heterogeneous
speeds ($+6.5\%$), and $2.82$ versus $3.11$ under the obstacle
penalty ($+9.5\%$). The margin is sensitive to the metric but the
sign is not. Under the Manhattan metric the ratio of longest to
shortest possible distance between two fixed points is compressed
($d_{\mathrm{Eucl}} \leq d_1 \leq \sqrt{2}\,d_{\mathrm{Eucl}}$),
which flattens the cost landscape and narrows the gap between
$p/d$ and $p/d^2$; under the obstacle variant, detour penalties
sharpen the cost landscape by punishing a handful of long crossings,
and the margin grows accordingly. The heterogeneous-speed variant
sits between the two: speed noise is a second source of variance
that slightly dampens the M4* advantage without reversing it. In
every case the cost-foreclosure argument --- that the quadratic
penalty protects a finite energy budget by discouraging
over-commitment to distant high-prior sites --- goes through without
reference to the specific distance metric. That the empirical
ranking tracks the argument across three qualitatively different
perturbations is evidence that the $p/d^2$ recommendation applies
to realistic deployments in which distance is not Euclidean.
